\documentclass{article}
\usepackage[T1]{fontenc}
\usepackage{lmodern}
\usepackage{graphicx}
\usepackage{amsmath,amssymb}
\usepackage[a4paper,margin=2cm]{geometry}
\usepackage[absolute,overlay]{textpos}
\setlength{\TPHorizModule}{1mm}
\setlength{\TPVertModule}{1mm}
\usepackage[indonesian]{babel}
\usepackage{hyperref}
\setlength{\emergencystretch}{2em}
% unit: Halaman 49

\begin{document}

% === Halaman 49 (llm) ===

\textbf{Contoh:} Misalkan $y^2 \equiv x^3 - x + 188 \pmod{751}$. Kurva eliptik ini memiliki $N = 727$ buah titik.

\begin{itemize}
    \item Misalkan karakter yang akan dikodekan adalah huruf 'B', yang dikodekan menjadi nilai 11.
    \item Pilih $k = 20$, maka $x = mk + 1 = (11)(20) + 1 = 221$. Sulihkan $x = 221$ ke dalam kurva eliptik $y^2 \equiv x^3 - x + 188 \pmod{751} \equiv 456 \pmod{751}$. Tidak ada nilai $y$ yang memenuhi.
    \item Coba untuk $x = mk + 2 = (11)(20) + 2 = 222$. Sulihkan $x = 222$ ke dalam kurva eliptik $y^2 \equiv x^3 - x + 188 \pmod{751}$. Juga tidak ada nilai $y$ yang memenuhi.
    \item Coba untuk $x = mk + 3 = (11)(20) + 3 = 223$. Sulihkan $x = 223$ ke dalam kurva eliptik $y^2 \equiv x^3 - x + 188 \pmod{751}$. Juga tidak ada nilai $y$ yang memenuhi.
    \item Coba untuk $x = mk + 4 = (11)(20) + 4 = 224$. Sulihkan $x = 224$ ke dalam kurva eliptik $y^2 \equiv x^3 - x + 188 \pmod{751}$. Diperoleh $y = 248$. Jadi, karakter 'B' dikodekan menjadi titik $(224, 248)$ pada kurva eliptik.
    \item Pada proses \textit{decoding}, hitung $m = \lfloor (x - 1)/k \rfloor = \lfloor (224 - 1)/20 \rfloor = \lfloor 11.15 \rfloor = 11$. Jadi pesan semula adalah huruf 'B'.
\end{itemize}

\end{document}
